\documentclass[11pt]{article}
\usepackage[a4paper,margin=26mm]{geometry}
\usepackage{amsmath,amssymb,mathtools}
\usepackage{unicode-math}
\setmainfont{Latin Modern Roman}
\setmathfont{Latin Modern Math}
\usepackage{xurl,hyperref}
\hypersetup{hidelinks,pdftitle={UnsolvedMath Problem 10400015}}
\newcommand{\sref}[2]{\hyperlink{source:#1:#2}{[S#2]}}
\newcommand{\cataloguescope}[1]{#1}
\begin{document}
% UnsolvedMath ID 10400015; source catalogue code AMR-103-0015
\setcounter{section}{10400014}
\section{Problem 1.15 — (M.}\label{10400015}
{\small\sffamily UnsolvedMath ID 10400015 · Topology}
\subsection{Definitions and mathematical statement}

\paragraph{Shared notation.}$\mathbb N=\{1,2,\ldots\}$, and $\mathbb Z,\mathbb Q,\mathbb R,\mathbb C$ denote the integers, rationals, reals and complex numbers. Any other conventions are specified locally.


A knot or link is a smooth oriented embedding of one or more circles in the oriented $3$-sphere, considered up to orientation-preserving ambient isotopy. Write $O$ for the unknot, which bounds an embedded disk. The HOMFLY (skein) polynomial $P_L(l,m)$ is normalized by $P_O=1$ and $l^{-1}P_{L_+}-lP_{L_-}=mP_{L_0}$, where $L_+,L_-,L_0$ agree outside a ball and inside it have a positive crossing, a negative crossing and the oriented smoothing, respectively. The Kauffman polynomial is $F_L(a,z)=a^{-w(D)}[D]$, where $w(D)$ is diagram writhe. The regular-isotopy bracket has $[O]=1$, $[D_+]+[D_-]=z([D_0]+[D_\infty])$, and a positive or negative curl multiplies it by $a$ or $a^{-1}$ respectively; the two smoothings in this relation are unoriented. Put $Q_L(z)=F_L(1,z)$. For a nonzero Laurent polynomial, $\deg_x$ and $\min\deg_x$ are the greatest and least exponents of $x$ with nonzero coefficient, and $\operatorname{span}_x=\deg_x-\min\deg_x$. A Whitehead double $W_K$ is the satellite knot obtained from the two oppositely oriented parallel strands of $K$ joined by a Whitehead clasp; the source allows the clasp sign and twist choice of a Whitehead double.
\paragraph{Statement recorded in the supplied source.}
For every nontrivial knot $K$ and every Whitehead double $W_K$, is
\[\deg_mP_{W_K}(l,m)=2\deg_zF_K(a,z)+2?\]
Here these are greatest exponents in the indicated variables. The problem concerns general nontrivial companions and their Whitehead doubles, not just the untwisted doubles of prime alternating knots discussed in the source update. \sref{10400015}{2}
\subsection{Short English statement}
Is the Alexander-variable degree of a Whitehead-double HOMFLY polynomial always twice the companion Kauffman z-degree plus two?
\subsection{Sources}
\begin{description}
\item[\hypertarget{source:10400015:1}{S1}] ULAM. AI and the UnsolvedMath Contributors, \emph{UnsolvedMath: A Curated Collection of Open Mathematics Problems}, record 10400015 (AMR-103-0015). CC BY 4.0. \url{https://huggingface.co/datasets/ulamai/UnsolvedMath}.
\item[\hypertarget{source:10400015:2}{S2}] T. Ohtsuki (editor), \emph{Problems on invariants of knots and 3-manifolds}, Geometry \& Topology Monographs 4 (2002), pp. 377--572 (published 2004), item 1.15 and its surrounding definitions and remarks. \url{https://arxiv.org/abs/math/0406190v1}.
\end{description}
\label{end:10400015}
\end{document}
